% Results (Gate B): seven sections, one per main figure; each heading is that figure's title (check_tex.py checks it). Every
% number is a ledger value (\val), every main panel is called out, procedures are in the past tense and findings in the
% present. In the preprint each section is followed by its figure and legend (\resultsfigure).

\subsection{A whole-CNS connectome learns to plan and to remember}

We built a spiking model of the complete central nervous system (CNS) of a male fly from its connectome\cellcite{berg2026}:
\val{L01} leaky integrate-and-fire neurons with spike-frequency adaptation, joined by \val{L02} million fixed connections
(\figs{1A}{1B}). Training could change \val{L05} gains and thresholds, set per cell type, and the \val{L06} weights of a linear
readout, but never the wiring (\fig{1E}). The model sees each task through \val{L03} visual input cells, and \val{L08} of its
\val{L04} readout cells (\val{L07}) receive a synapse directly from them (\fig{1E}). Training reached the whole CNS: it changed
the output gains of nearly every cell, in every division (\figr{S1A}{S1C}).

We trained the model on two tasks (\fig{1F}). To plan, it pushes boxes onto targets in a ladder of \val{L09} maps of rising
difficulty, each with a certified optimal solution. To remember, it is shown which of two targets is cued, holds the cue through
a blank delay of up to \val{H12}~ms and then, from the board without the target, pushes the box toward the cued one. During a
memory trial, cells throughout the CNS spike in the cue and the input, and few spike in the delay (\figs{1C}{1D}).

The model learns both tasks reproducibly. Learning cost tracked the certified difficulty of the maps in all \val{L10} replicate
sets (Spearman ρ = \val{L11} to \val{L12}; \fig{1G}; \figs{S1D}{S1E}). The final brain with the final readout solves all
\val{L13} maps, whereas no other pairing of brain and readout checkpoints solves more than \val{L14} (\fig{1H}), so brain and
readout learned together; \video{S1} shows the hardest map solved in its \val{L18} certified moves. In the memory task, training
used \val{A01} of \val{A04} combinations of delay and input timing. We scored each timing by net memory use: \val{A02} when the
model solves the histories of both cues, 0 when it solves only one, as a model that ignores the cue would, and $-$\val{A02} when
it solves neither. We trained four learners: the readout alone or the brain and readout in scene 1, and the readout alone
in scenes 2 and 3, two further board layouts of the task. In all \val{L15} replicates of the four learners, the model uses the
cue at the trained timings and not at the others (Holm-adjusted \Pv{} = \val{L16} to \val{L17}; \fig{1I}), and every run
replays bit for bit (\fig{S1I}). We next asked how a network without a dedicated memory circuit holds the cue through the
silent delay.
\resultsfigure{1}

\subsection{Working memory is held in an activity-silent intrinsic trace}

To find where the cue is held, we recorded every state variable of every cell through single trials (\fig{2A}; \video{S2}).
Spikes that tell the two cues apart stop \val{H01} to \val{H02}~ms after the cue ends in the four learners (\fig{2B};
\fig{S2A}), yet the cue persists. In the trial of \fig{2A}, \val{H03} cells hold a difference in spike-frequency adaptation
between the cues through the delay (\fig{2C}; \fig{S2C}). Of all state variables, only adaptation keeps the cue, and its
difference decays with the adaptation time constant of \val{H04} to \val{H05}~ms (\fig{2D}; \fig{S2D}). Holding the trace costs
no spikes: in none of the four learners does a cell spike differently during the last \val{H11}~ms of the \val{H12}-ms delay
(\fig{2I}; \fig{S2H}).

The next input reads the trace out. At the end of the delay the cue difference in membrane voltage comes from adaptation, but at
the decision it comes from synaptic input (\fig{2E}), and at least \val{H06} of the decision-time drive is synaptic at all
\val{H07} timings (\fig{2H}; \fig{S2E}). Input also renews the trace, which otherwise decays as fitted (\fig{2J}).

The trace is sufficient and specific. Moving the full state at the end of the delay into another trial recreates the native
decision (cosine \val{H08}; \fig{2F}; \fig{S2F}). Swapping the states of the two cues flips the decisions at all trained timings
in all four learners (Holm-adjusted \Pv{} = \val{H09} to \val{H10}; \fig{2G}). The input reads the trace through the
cells that hold it: the \val{H14} of cell types that held the trace supply \val{H13} of the cue drive into readout types, in the
brain and the ventral nerve cord alike (\figs{2H}{2K}). The cue at the decision is spread over many cells in every learner
(\fig{S2B}). We next asked where in the CNS the trace sits, and whether any circuit is special.
\resultsfigure{2}

\subsection{The trace is distributed, and no circuit is special}

Memories are often assigned to ensembles of cells\cellcite{roy2022}. We took the \val{D15} of cells that carry the most cue
difference at the decision as the cue ensemble of each of two scenes (\val{D03} cells in scene 1; \fig{3A}; \fig{S3A}). The
ensemble is densely wired, with \val{D01} and \val{D02} times the internal connections of \val{D14} matched sets in the two
scenes (\fig{3B}), and it is made of the same cell groups as the rest of the CNS (\fig{S3B}).

Yet the ensemble is not necessary. Silencing it at the decision changes the decision no more than silencing matched sets
(Holm-adjusted \Pv{} = \val{D04}; \fig{3C}), and cutting its internal synapses does no more than cutting matched synapses
(\Pv{} = \val{D05}; \fig{3D}; \fig{S3C}). Whole populations are dispensable: with half or all of the Kenyon cells or the
antennal-lobe inhibitory neurons removed, the network solves the trained histories that the healthy network solves, \val{D06}
of \val{D07} (equivalence test, Holm-adjusted \Pv{} = \val{D08}; \fig{3E}; \fig{S3D}).

The cue concentrates where it is read, not in olfactory circuits: the readout cells carry \val{D09} of the cue difference,
\val{D10} to \val{D11} times their share of cells (\fig{3F}). None of this depends on how the connectome becomes a network. All
\val{D12} runs learn under \val{D13} sets of connection rules, and the signatures of healthy circuits hold under each (\fig{3G};
\fig{S3E}). If no circuit is special, the cells that hold the trace might still share a property, and we asked what marks them.
\resultsfigure{3}

\subsection{Trace-carrying cells have a molecular signature: \gene{SK}}

To ask what marks the cells that hold the trace, we joined published single-cell atlases of the fly brain and ventral nerve
cord\cellcite{davie2018,ozel2021,davis2020,lu2023,park2025,allen2020} to the connectome, by cell type and, in the ventral nerve cord, through
\val{M05} hemilineages (\figr{4A}{4C}; \figs{S4A}{S4B}). The join reproduces the atlas clusters (Spearman ρ = \val{M07}) and the
fast transmitter predicted from electron microscopy\cellcite{eckstein2024} in \val{M06} of \val{M05} hemilineages (\fig{4H};
\figr{S4C}{S4E}). It maps \val{M09} of the trace carriers (the cells that hold the trace at the first decision), up from \val{M08} by
cell type alone, and \val{M10} in the ventral
nerve cord (\fig{4F}).

As we declared before testing, the carriers are enriched for \gene{SK}, which encodes a small-conductance calcium-activated
potassium channel. Carriers exceed matched sets in \gene{SK} and in an adaptation index (adjusted \Pv{} ≤ \val{M01}), and
\gene{SK} exceeds a null that shuffles the mapping (adjusted \Pv{} = \val{M02} to \val{M03}; \fig{4D}; \fig{S4F}). The ventral
nerve cord, whose data never shaped the prediction, confirms it: across \val{M05} hemilineages, the share of carriers rises with
\gene{SK} (Spearman ρ = \val{M04}; \fig{4E}), against the background of the full channel panel (\fig{4G}).

The carriers are accessible in the fly. Four carrier cell types with split-GAL4 driver lines send \val{M11} to \val{M12} of
their cue drive into readout types (\fig{4I}), and the cue reaches the motor neurons through \gene{SK}-rich hemilineages
(\fig{S4G}). Adaptation that varies from cell to cell, as the atlases suggest, leaves the mechanism intact: brains with
atlas-informed heterogeneity keep the memory, its lock to the trained timings and the silent trace in all \val{M14} runs
(Holm-adjusted \Pv{} = \val{M15}; \fig{4J}). If the trace is the memory, something else must decide when it can be read.
\resultsfigure{4}

\subsection{Access to the memory is locked to the timing of learning}

The readout uses the trace only at the timings it was trained on. Net memory use is \val{A02} at all \val{A01} trained timings
in every replicate and averages \val{A03} at the other timings (Holm-adjusted \Pv{} = \val{A05}; \fig{5B}; \fig{S5B}), and the
lock holds in \val{A07} of \val{A08} runs across nine scenes (\fig{5C}).

The lock is not a loss of the trace. At untrained timings the trace reaching the decision is no smaller (median size \val{A11}
against \val{A12} at trained timings), but its push toward the correct target collapses (Holm-adjusted \Pv{} = \val{A13};
\fig{5D}; \fig{S5C}). After the untrained \val{G01}-ms delay, the decision-time pattern is nearly the mirror image of the
trained \val{G06}-ms pattern (cosine \val{A09}; \fig{5A}; \fig{S5A}), and \val{A10} of the strongest cells change sign
(\fig{5E}; \fig{S5D}). In the space of decisions, trained timings push toward the correct target, and untrained timings do not
(\fig{5F}).

The cause is a drift in how the trace is read. Injected at seven delays, one fixed trace pattern pushes toward the target only at
the trained delays (\val{A14} to \val{A15}; \fig{5G}; \fig{S5E}). The operator that turns the trace into a push therefore changes
with delay while the trace persists, and we call this operator drift. Every cell group turns with it (\fig{S5D}).
Training across jittered timings widens use inside the jittered range (\val{A16}) but not beyond it (\val{A17}; \fig{5H}; \fig{S5F}), and \video{S3} shows one trace read at a trained
and an untrained delay. We next asked whether access is also locked to the network's excitability.
\resultsfigure{5}

\subsection{The wiring is robust; access follows the trained states}

The connectome keeps the readout within reach. With \val{G02} of the other cells removed at random, \val{E01} of the readout
cells remain reachable from the visual input, while the cells on routes fall from \val{E02} to \val{E03} (\figs{6A}{6B};
\figs{S6A}{S6B}). Routes thin through every cell class, but none breaks: three-step routes fall to \val{E04} of intact, and
\val{E06} of \val{E05} cell types still carry them (\figs{6G}{6H}; \fig{S6G}).

The dynamics are stable. A single extra spike leaves nothing at the decision in \val{E07} of \val{E08} kicks (\fig{6C};
\fig{S6C}), and the mean firing rate stays between \val{E09} and \val{E10}~Hz for threshold changes of up to \val{G09}~mV
(\fig{6D}).

Access follows the states of training. Decoder theory predicts that a readout trained across excitability states keeps working
where one trained in a single state fails\cellcite{haimerl2025}, and the model bears this out. Training across threshold states
widens the access window \val{E11}-fold (\val{G11} interval \val{E12} to \val{E13}; Holm-adjusted \Pv{} = \val{E16}). It also
keeps memory under threshold jitter (\val{E14} against \val{E15}) and raises tolerance to cell loss \val{E17}-fold (\val{E18}
to \val{E19}; \fig{6E}; \figs{S6D}{S6E}). The mechanism is the same in both: a silent adaptation trace read by the input
(\fig{6F}). We then asked what happens when disease shifts the state after learning.
\resultsfigure{6}

\subsection{Disease steals access, not memories}

We screened \val{Q01} fly pathologies for credible operators, requiring published evidence, a mechanism the model can express,
healthy circuits that match the fly and the fly phenotype reproduced in two brains (\fig{7A}; \fig{S7B}; \tab{S1}). Of these,
\val{Q02} passed, and only one reaches visual memory: the pan-neuronal loss of the A-type potassium channel Kv4 that Aβ42 causes
in flies\cellcite{ping2015} (\fig{7F}).

A lesion after learning removes access, whereas learning under it is spared (\fig{7B}). At dose 1, the strength that reproduces the
fly phenotype, the Kv4-loss lesion lowers the trained histories that the learned readout solves from \val{Q04} to \val{Q03},
and memory falls with dose (Holm-adjusted \Pv{} = \val{Q05}; \fig{7C}; \fig{S7C}). Readouts trained in the lesioned brain solve \val{Q07} of \val{Q04} in every one of
\val{Q06} runs (\fig{7C}; \fig{S7D}). Retraining the readout after the lesion recovers \val{Q08} of the loss (\fig{7D};
\fig{S7E}), but restoring excitability late, after learning under the lesion, removes access again and leaves \val{Q09} of
\val{Q04} (BH-adjusted \Pv{} ≤ \val{Q10}; \fig{7E}; \fig{S7F}). Memories learned under one lesion pattern are likewise lost
under another (\fig{7I}).

An atlas of Aβ42 flies constrains how the lesion acts: neuronal \gene{Shal} transcripts rise in it (\val{Q11}
log\textsubscript{2})\cellcite{park2025}, so the Kv4 loss must arise after transcription (\fig{7G}). Patterned by \gene{Shal}, the lesion drives
\val{Q12} to \val{Q13} of cells above \val{G03}~Hz at \val{G04} to \val{G05} times the dose (\fig{7H}; \figr{S7H}{S7J}).

To ask whether the lesion damages the stored trace or its reading, we moved the full state of each trial at the end of the delay
between the healthy and the lesioned brain (\fig{7J}; \video{S4}). Recall needs the learning state at both ends. Readouts solve
all \val{K01} trained histories when the trace is written and read in the state of learning, and at most \val{K02} when either
end differs. At dose 1, a healthy trace read under the lesion gives \val{K03} of \val{Q04} and a lesion-written trace read in
health \val{K04}; for readouts trained under the lesion, these are \val{K05} and \val{K06}.

The lesion changes both ends. It writes the trace into \val{K11} to \val{K12} cells instead of \val{H03}, nearly orthogonal to
the healthy trace (cosine \val{K07} to \val{K08}). The hold stays silent only at the lowest dose: \val{K13} cells spike
differently late in the delay at dose \val{G10}, against \val{K14} at dose 1. Read under the lesion, the same healthy
trace evokes a decision-time pattern nearly orthogonal to its healthy reading (median cosine \val{K09} to \val{K10}). The
lesioned brain also spends more for less: across doses, a trial costs \val{Q21} to \val{Q22} times the spikes of a healthy trial for \val{Q23}
to \val{Q24} bits about the cue, against \val{Q20} bits in health (\fig{S7I}). An aging operator built from transcripts of the
aging fly\cellcite{lu2023} dims Kenyon-cell odor responses with its strength (\fig{S7K}). In the model, memory failure is a
change of state at writing and at reading, not a loss of storage.
\resultsfigure{7}
